%% BIBM 2026 — IEEE conference format.
%% Compile on Overleaf with the IEEE BIBM template (\documentclass[conference]{IEEEtran}).
%% HONESTY CONTRACT (do not overclaim — see notes in each section):
%%   - spinopelvic params: implemented + geometry unit-tested; manual-measurement
%%     validation is ONGOING (not done). PI is posture-invariant; SS/PT/LL on supine CT
%%     are surrogates. SVA/TPA are out of scope (FOV/posture) — say so.
%%   - DRR/XR: projection + label transfer WORK; the full 1k release is forthcoming;
%%     the XR level-localizer is a FRAMEWORK (not trained).
%%   - viewer: live single-case web demo; engine-swap contract is DESIGNED (not dual-engine).
%%   - post-op synthesis: prototype (illustrative), not validated.
%% \TODO{...} marks a number/figure to drop in after running tools/cohort_stats.py.
\documentclass[conference]{IEEEtran}
\usepackage{cite,amsmath,graphicx,booktabs,url,xcolor}
\newcommand{\TODO}[1]{\textcolor{red}{[TODO: #1]}}
\newcommand{\dataset}{CTSpinoPelvic1K}
\newcommand{\xrdataset}{XRSpinoPelvic1K}
\newcommand{\ostk}{OpenSpineToolkit}

\begin{document}
\title{From Segmented CT to Spinopelvic Measurement:\\
An Open Geometric Engine, a CT$\rightarrow$Radiograph Bridge,\\ and a Deployable Viewer}

\author{\IEEEauthorblockN{Gregory Schwing\,$^{\dagger}$ and the OpenSpine Consortium}
\IEEEauthorblockA{$^{\dagger}$Correspondence: [personal e-mail struck on copy, 2026-10-08]}}

\maketitle

\begin{abstract}
Spinopelvic parameters---pelvic incidence (PI), sacral slope (SS), pelvic tilt (PT)
and lumbar lordosis (LL)---drive surgical planning for adult spinal deformity, yet they
are still measured by hand on radiographs, with substantial inter-rater variability, and
they are rarely available from the cross-sectional CT that is increasingly acquired first.
We present an open, reproducible pipeline that turns a densely segmented CT volume into
spinopelvic measurements and brings the same measurements to the radiographic domain.
First, a geometric engine (\ostk) derives PI/SS/PT/LL directly from world-space anatomical
landmarks fit to the segmentation (robust femoral-head sphere fit, sacral-endplate and
vertebral-endplate plane fits), emitting a per-case JSON contract with explicit
quality-control flags; the construction is internally consistent ($|\mathrm{SS}+\mathrm{PT}-\mathrm{PI}|<1^\circ$).
Second, we close the CT--radiograph gap two ways: an analytic 2D port computes the
identical angles on lateral radiographs (cross-validated against the 3D engine on
identical geometry), and a digitally-reconstructed-radiograph (DRR) pipeline projects the
CT \emph{and its dense labels} through one shared geometry to produce paired synthetic
radiographs, per-pixel label maps and per-level landmarks with no manual annotation
(\xrdataset). Third, a browser viewer renders a case and animates each angle construction
from the JSON contract, which is designed to decouple the measurement engine from the
interface so a learned engine can later replace the geometric one. We report the parameter
distribution over the \dataset{} cohort, the 2D--3D equivalence, and an illustrative
post-operative alignment-simulation mode. All code is released under Apache-2.0.
\end{abstract}

\begin{IEEEkeywords}
spinopelvic parameters, CT segmentation, digitally reconstructed radiograph, surgical
planning, open-source, biomedical image analysis
\end{IEEEkeywords}

\section{Introduction}
Sagittal spinopelvic alignment governs outcomes in adult spinal deformity surgery.
The core parameters---pelvic incidence (PI), sacral slope (SS), pelvic tilt (PT) and
lumbar lordosis (LL)---and the PI$-$LL mismatch underlie widely used classification and
realignment schemes \cite{schwab2012,lafage2016}. In practice these are measured manually
on standing lateral radiographs, a process that is time-consuming and shows non-trivial
inter- and intra-rater variability. At the same time, cross-sectional CT is frequently the
first study available (trauma, oncologic, pre-operative), and modern automated
segmentation now labels the spine and pelvis on CT at scale
\cite{verse,totalsegmentator}. This raises a practical question: can spinopelvic
assessment be derived \emph{automatically and reproducibly} from a segmented CT, and can
the result be made usable in the radiographic workflow clinicians already trust?

We answer this with an open, end-to-end pipeline built on a densely labeled CT resource,
\dataset{}\TODO{cite the descriptor once arXiv/MLHC id is known}, and make three
contributions:
\begin{enumerate}
\item \textbf{A geometric measurement engine} (\ostk) that computes PI/SS/PT/LL from a CT
segmentation using world-space landmark geometry, with a machine-readable output contract
and per-measurement quality-control (QC) flags. The construction is internally consistent
by design and unit-tested.
\item \textbf{A CT$\rightarrow$radiograph bridge}: an analytic 2D port that reproduces the
identical angles on lateral radiographs (cross-validated against the 3D engine), together
with a DRR pipeline that projects the CT and its dense labels through one shared geometry
to yield paired synthetic radiographs, label maps and per-level landmarks
(\xrdataset)---annotation-free supervision for radiograph-based analysis.
\item \textbf{A deployable, open viewer} that renders a case and animates each angle from a
JSON contract designed to separate the measurement engine from the interface, plus an
illustrative post-operative alignment-simulation mode.
\end{enumerate}
Everything is released under Apache-2.0 for reproducibility. We are explicit about scope:
PI is posture-invariant and therefore directly meaningful on supine CT, whereas SS, PT and
LL on supine CT are reported as surrogates; SVA and the T1-pelvic angle are out of scope on
this acquisition (Section~\ref{sec:disc}).

\section{Related Work}
Automated vertebral labeling and spine/pelvis segmentation on CT are now mature
\cite{verse,totalsegmentator,ctspine1k,ctpelvic1k}. Prior automation of spinopelvic
parameters has largely targeted radiographs via landmark regression, or proprietary
vendor tools; reproducible, openly licensed pipelines that operate from CT segmentation
and expose an auditable geometric construction are scarce. DRR synthesis is well
established for 2D/3D registration and for training radiograph models
\cite{drr}; we use it here specifically to transfer \emph{dense anatomical labels} from CT
into the radiographic domain so that radiograph measurements inherit CT-derived geometry.
Our contribution is integrative and open: a single, inspectable path from CT labels to
spinopelvic measurement, to radiograph transfer, to a deployable viewer.

\section{Methods}
\subsection{Dense CT labels (substrate)}
The pipeline consumes a CT segmentation in a fixed label scheme covering the lumbar
vertebrae, sacrum and S1, the hips and femora, and (FOV-permitting) the lower thoracic
vertebrae and ribs. These labels are produced by the \dataset{} pipeline
\TODO{1-2 sentences + cite descriptor}; this paper concerns everything \emph{downstream}
of the labels. All geometry is computed in world millimetres from the NIfTI affine.

\subsection{Geometric spinopelvic engine (\ostk)}
\label{sec:ostk}
\textbf{Landmarks.} Femoral-head centres are estimated by a robust two-stage sphere fit
seeded from the acetabular interface and refined by shell-outlier rejection; the
bicoxofemoral axis (line through the two centres) defines the patient left--right
direction and hence the sagittal plane used by every angle. The S1 superior endplate and
each vertebral endplate are fit by a total-least-squares plane to an anterior band of the
endplate surface, which rejects posterior elements and osteophytes.

\textbf{Parameters.} With the S1 endplate normal $\mathbf{n}_{S1}$, the femoral-head
midpoint $\mathbf{m}$ and the endplate centroid $\mathbf{c}$, all projected into the
sagittal plane:
PI is the angle between $\mathbf{n}_{S1}$ and the radius $\mathbf{m}\!\to\!\mathbf{c}$;
SS is the angle of $\mathbf{n}_{S1}$ to vertical; PT is the angle of
$\mathbf{m}\!\to\!\mathbf{c}$ to vertical. LL is the signed Cobb angle between the L1 and
S1 superior endplates, decomposed into per-level segments (L1--L2 $\ldots$ L5--S1). The
construction enforces the morphological identity $\mathrm{PI}=\mathrm{SS}+\mathrm{PT}$,
checked per case ($|\mathrm{SS}+\mathrm{PT}-\mathrm{PI}|<1^\circ$).

\textbf{Output contract and QC.} Each measurement is emitted as JSON with the value, the
world-mm landmarks, plane-fit residuals, a method version, and QC flags
(e.g.\ \texttt{low\_voxels}, high fit residual, identity violation) so that unreliable
cases are surfaced rather than silently returned. PI$-$LL mismatch and SRS--Schwab
severity modifiers are computed from the parameters \cite{schwab2012}.

\subsection{CT$\rightarrow$radiograph bridge}
\label{sec:bridge}
\textbf{Analytic 2D port.} The same angle definitions are implemented for 2D lateral
geometry; on identical input geometry the 2D and 3D engines agree to numerical tolerance
(Section~\ref{sec:results}), so a measurement made on a lateral radiograph is the same
quantity the CT engine computes.

\textbf{DRR + label projection (\xrdataset).} We render synthetic radiographs by
attenuation integration (parallel-beam ray casting) along the projection axis, with a
bone-emphasis mode and correct voxel spacing, in lateral and AP views. Crucially, the
dense \emph{labels} are projected through the \emph{same} projection geometry, producing,
per case and view: a synthetic radiograph, a per-pixel projected label map, and per-level
landmarks (centroid and bounding box). Because image and labels share one geometry, the
2D supervision is exact and requires no manual radiograph annotation. The builder runs
over \dataset{} CT--label pairs; we show a fully rendered example
(Fig.~\ref{fig:drr}) and report the projection alignment. A heatmap-based level-localizer
\emph{framework} (per-level Gaussian targets, a nearest-level baseline, and a compact
U-Net regressor) is provided; training and evaluation on the rendered set are future work.

\subsection{Deployable viewer and engine contract}
A browser viewer (NiiVue) loads a compact per-case bundle (cropped, bone-masked,
downsampled CT $\sim$1\,MB; the label map; and a \texttt{metrics.json}). Clicking a
parameter animates its geometric construction---solid anatomical lines, dashed
constructions and the angle arc---mapped from world mm to the screen each frame, so the
overlay stays correct as the user scrolls. The \texttt{metrics.json} schema stores the
construction in world coordinates independently of the renderer; it is designed so that a
learned engine can later emit the same contract and replace the geometric one without
changing the interface. The viewer is an educational, browser-only demonstration (NIfTI,
not a clinical PACS) and is hosted publicly.

\subsection{Post-operative alignment simulation (prototype)}
For planning visualisation, \ostk{} can simulate a correction: a mobile spinal segment is
rotated about a technique-specific fulcrum (interbody/osteotomy), interbody cages are
rendered, the pelvis is analytically re-oriented to a target PT, and a smooth displacement
field warps the CT to a plausible post-operative image; age-adjusted realignment targets
follow \cite{lafage2016}. This mode is an illustrative prototype and is not validated
against post-operative imaging.

\section{Results}
\label{sec:results}
\textbf{Internal consistency.} Across the cohort the PI$=$SS$+$PT identity holds to
$<1^\circ$ on \TODO{N} cases that pass QC; \TODO{report \% with each QC flag} were flagged
for low voxel support or high fit residual and excluded from summary statistics.

\textbf{2D--3D equivalence.} On matched geometry the analytic 2D port reproduces the 3D
PI/SS/PT/LL to within \TODO{$\epsilon$}$^\circ$ (unit-tested), confirming that radiograph
and CT measurements are the same quantity.

\textbf{Cohort parameter distribution.} Running the engine over the \dataset{} cohort
yields the distributions in Table~\ref{tab:dist} (mean$\pm$SD; valid-measurement rate per
parameter). \TODO{fill Table~\ref{tab:dist} from tools/cohort\_stats.py output}

\begin{table}[t]
\caption{Spinopelvic parameters over the \dataset{} cohort (supine CT). SS/PT/LL are
supine surrogates; PI is posture-invariant.}
\label{tab:dist}
\centering
\begin{tabular}{lccc}
\toprule
Parameter & Mean$\pm$SD ($^\circ$) & Range & Valid rate \\
\midrule
PI & \TODO{} & \TODO{} & \TODO{} \\
SS & \TODO{} & \TODO{} & \TODO{} \\
PT & \TODO{} & \TODO{} & \TODO{} \\
LL & \TODO{} & \TODO{} & \TODO{} \\
\bottomrule
\end{tabular}
\end{table}

\textbf{CT$\rightarrow$radiograph.} Fig.~\ref{fig:drr} shows a rendered lateral DRR with
its auto-projected per-level labels; projected bone footprints align with DRR bone
intensity \TODO{quantify overlap on the example/subset}.

\textbf{Viewer.} Fig.~\ref{fig:viewer}\TODO{screenshot} shows the live viewer with an
animated PI construction and the report panel.

\begin{figure}[t]\centering
%\includegraphics[width=\columnwidth]{fig_drr_example.png}
\fbox{\parbox{0.9\columnwidth}{\centering\vspace{2.2cm}\TODO{Fig: lateral DRR + projected
level labels (example case)}\vspace{2.2cm}}}
\caption{Synthetic lateral radiograph (DRR) with dense labels projected through the same
geometry: paired image, label map and per-level landmarks, annotation-free.}
\label{fig:drr}
\end{figure}

\begin{figure}[t]\centering
\fbox{\parbox{0.9\columnwidth}{\centering\vspace{2.2cm}\TODO{Fig: viewer screenshot ---
animated PI/SS/PT/LL construction + report}\vspace{2.2cm}}}
\caption{Browser viewer: angle constructions animated from the \texttt{metrics.json}
contract; engine is decoupled from the interface.}
\label{fig:viewer}
\end{figure}

\section{Discussion and Limitations}
\label{sec:disc}
\textbf{Scope.} PI is a morphological constant and is therefore directly valid on supine
CT; SS, PT and LL depend on posture, so their absolute values on supine CT are reported as
surrogates rather than standing equivalents. SVA and the T1-pelvic angle require a C7/T1
reference and standing acquisition and are out of scope on this data
(C7 is outside the field of view); the engine marks them as such rather than reporting
unreliable values.

\textbf{Validation.} The geometric construction is unit-tested for correctness (identity,
plane/sphere fits, 2D--3D equivalence), but agreement against manual radiographic
measurement (MAE, ICC, Bland--Altman) is ongoing and is the primary next step.

\textbf{Maturity.} The radiograph level-localizer is provided as a framework, not a trained
model; the \xrdataset{} full release is forthcoming (the generation pipeline is complete and
demonstrated); the viewer is a single-case educational demonstration with a contract
designed---but not yet wired---for a dual (geometric/learned) engine; and the
post-operative mode is an illustrative prototype. We state these explicitly so the
contribution is not over-read.

\textbf{Why it matters.} Even within this scope, the pipeline gives an open, auditable,
reproducible route from routinely acquired CT to spinopelvic geometry, and---through exact
label transfer---an annotation-free path to bring the same measurements to radiographs.

\section{Conclusion}
We presented an open pipeline from segmented CT to spinopelvic measurement: a geometric
engine with an auditable contract, a CT$\rightarrow$radiograph bridge (exact 2D port plus
label-carrying DRRs), and a deployable viewer. The components are released under Apache-2.0;
manual-measurement validation and a trained radiograph localizer are the immediate next
steps.

\section*{Reproducibility}
Code: \ostk{} (\url{https://github.com/Gregory-Schwing-MD-PhD/OpenSpineToolkit}); DRR/\xrdataset{}
pipeline and the viewer are open under Apache-2.0\TODO{add repo URLs}. Demo:
\url{https://openspineconsortium.com/pacs/}.

\begin{thebibliography}{1}
\bibitem{schwab2012} F. Schwab \emph{et al.}, ``Scoliosis Research Society--Schwab adult
spinal deformity classification: a validation study,'' \emph{Spine}, 2012.
\bibitem{lafage2016} R. Lafage \emph{et al.}, ``Age-adjusted alignment goals in adult
spinal deformity,'' \emph{Spine}, 2016.
\bibitem{verse} A. Sekuboyina \emph{et al.}, ``VerSe: a vertebrae labelling and
segmentation benchmark,'' \emph{Medical Image Analysis}, 2021.
\bibitem{totalsegmentator} J. Wasserthal \emph{et al.}, ``TotalSegmentator: robust
segmentation of 104 anatomic structures in CT images,'' \emph{Radiology: AI}, 2023.
\bibitem{ctspine1k} Y. Deng \emph{et al.}, ``CTSpine1K: a large-scale dataset for spinal
vertebrae segmentation in CT,'' 2021.
\bibitem{ctpelvic1k} P. Liu \emph{et al.}, ``CTPelvic1K: pelvic bone segmentation in CT,''
\emph{IJCARS}, 2021.
\bibitem{drr} \TODO{DRR reference, e.g. DiffDRR/DeepDRR or a classic Siddon ray-casting
citation}
\bibitem{nnunet} F. Isensee \emph{et al.}, ``nnU-Net: a self-configuring method for deep
learning-based biomedical image segmentation,'' \emph{Nature Methods}, 2021.
\end{thebibliography}
\end{document}
